\section{Primary construction: removing unsafe data during a finite prefix}
\subsection{The command language and the unsafe first root}
Let $P$ be Cantor pairing and $I(a,b)=4P(a,b)$. Write the retained tuple as $(t,w,d,h)$. Order 201 means
\[
\mathsf{NIP}:0,\qquad \mathsf{MP}:1,\qquad
\mathsf{DROP}:2,\qquad \mathsf{AX}(c;x,y):c+3.
\]
The non-AX field lists are $[0]$, $[1,x]$, and $[2]$ respectively; AX fields are $[c+3,y,x]$. The primitive bodies and their all-natural command meanings are inherited. Independent control resolvers check the selector partitions, field arities, and normal returns for this exact source.

NIP tests $h=1$ \emph{before} popping. If it does not halt and $(x,v)=P^{-1}(w)$, it returns $(t,v,d,x)$. DROP returns $(h,P(t,w),d,h)$. MP$(x)$ puts $y=d+x$ and returns
\[
\begin{cases}(y,P(0,w),y,0),&I(t,y)=h,\\(0,w,y,0),&I(t,y)\ne h.\end{cases}
\]
AX puts $U=d,A=d+x,B=h+y,C=c$, returns the gated bundle as top, resets $d=h=0$, and pushes the seven logical rows, the dominated row, $Q_Uh$, and the old top. The bound $U\le A$ is automatic for every natural input.

Exact blank execution first reaches a complete zero-address source frame after 118,411 transitions with
\[
(N,T,W,A,D,H,X,Y,L)=(2,0,2,0,0,0,0,0,0).
\]
This is \emph{not} a theorem-stack seed: $2=P(0,1)$ contains truth above an unsafe tail. No consistency invariant is applied at that point.

\subsection{Four commands that discharge the unsafe suffix}
The exact candidate packing determines the next commands, without assumptions on ZFC. Since $2=P(0,1)$ and subsequent positive remainders are decremented before decoding, candidate 2 executes NIP,NIP. Candidate $3=P(2,0)$ executes DROP. Candidate $4=P(1,1)$ executes MP$(1)$ because $1=P(1,0)$ supplies the payload.

\begin{center}\small
\begin{tabular}{lrrrr}\toprule
Boundary & $T$&$W$&$D$&$H$\\\midrule
Initial complete frame&0&2&0&0\\
After first NIP&0&1&0&0\\
After second NIP&0&0&0&1\\
After DROP&1&0&0&1\\
After MP$(1)$&0&0&1&0\\\bottomrule
\end{tabular}
\end{center}
Neither NIP halts: both test cache zero, and the second only \emph{then} loads falsity. DROP has no halt test. The final MP fails because $I(1,1)=16\ne1$, clearing both active formula registers. It leaves only the harmless additive counter value 1.

An independently implemented interval-control interpreter and the arithmetic control resolver agree on all source edges. Both reproduce the 435-primitive finite segment, with root visits after 0, 63, 89, 166, and 435 primitives. No halting primitive occurs. The terminal root is
\[
                         (5,0,0,0,1,0,0,0,0).
\]
A counterfactual theoremhood assertion during the earlier segment would be wrong. A verified finite nonhalting prefix needs no such invariant; the invariant begins only at its safe endpoint.

\subsection{A complete literal binary prefix}
The proof does not rely solely on the source interpreter for initialization. Independent native and canonical simulation replays all 570,263 binary transitions from blank tape, never halting, and reaches native state 84, head $-9$, gap 131, with the displayed safe tuple and 176 natural spare words.

The full sparse tape is reconstructed and compared. Its minimum one-cell is $-10=h-1$, and its next one is $186=h+64+131$. Hence all 64 source-address cells are zero. This is a genuine root, not an interior dispatcher frame. Every primary word, every mixed spare value, and the entire zero exterior are checked.

\subsection{Strong proof compilation and arbitrary initial counters}
The inherited strong compiler supplies, for each numerical theorem $F$, a finite named-command block $\mathcal B_F$ which from every $(t,w,0,0)$ either halts or returns $(F,w,0,0)$. It quantifies over arbitrary incoming top and stack. Command renumbering transports it to order 201 by changing wire tags, not arithmetic meanings.

For clarity, its principal blocks can be given explicitly. Put
\[
\mathsf{Discard}=\mathsf{MP}(1);\mathsf{AX}(1;0,0);\mathsf{NIP}^{10},
\quad J(u)=\begin{cases}[],&u=0,\\{}[\mathsf{MP}(u)],&u>0.\end{cases}
\]
From $D=H=0$, Discard returns truth with the old stack and zero counter/cache. From a zero top/counter/cache, $J(u)$ sets the counter to $u$; omitting MP when $u=0$ avoids its zero-alias successful push. Truth uses Discard. A retained row $R_j(a,b,c)$ uses
\[
\mathsf{Discard};\mathsf{AX}(c;a,b);\mathsf{NIP}^j;
\mathsf{DROP};\mathsf{NIP}^{11-j}.
\]
The dominated row with $u\le a$ uses
\[
\mathsf{Discard};J(u);\mathsf{AX}(c;a-u,b);
\mathsf{NIP}^8;\mathsf{DROP};\mathsf{NIP}^{3}.
\]
The bundle uses $\mathsf{Discard};\mathsf{AX}(c;a,b);\mathsf{NIP}^{10}$ when its gates hold. All generated row entries are nonfalsity implication codes and the generalized cache has tag 2, so these selection suffixes restore the underlying stack exactly.

For generalization use
\[
\mathcal B_X;\mathsf{DROP};J(i);\mathsf{NIP};\mathsf{AX}(i;0,0);
\mathsf{NIP}^{9};\mathsf{DROP};\mathsf{NIP};\mathsf{NIP}.
\]
For modus ponens from $X$ and $I(X,Y)$ use
\begin{multline*}
\mathcal B_X;\mathsf{DROP};\mathcal B_{I(X,Y)};
\mathsf{NIP};\mathsf{DROP};\mathsf{NIP};\mathsf{MP}(Y);\mathsf{AX}(Y;0,0);\\
\mathsf{NIP}^{10};\mathsf{DROP};\mathsf{NIP};\mathsf{NIP}.
\end{multline*}
The AX reset has equal selector/counter parameter and therefore zero gate; the last selection restores $Y$ and the old stack, or halts earlier on falsity. These exact contracts, including the zero implication alias, prove the compiler by induction on a finite numerical derivation.

\begin{theorem}[Safe root with arbitrary counter]
At a genuine root with $N=n,D=d$ arbitrary naturals, $A=X=Y=L=0$, and theorem-valued top, cache, and stack heads, source halting is equivalent to numerical inconsistency.
\end{theorem}
\begin{proof}
Soundness permits every $d$: AX uses $U=d\le d+x=A$, and all other commands retain their theoremhood arguments. Only cached falsity can trigger halt.

For completeness, append DROP,NIP,NIP to $\mathcal B_1$. Prefix by $\mathsf{AX}(0;1,y)$ with $y\ge1$. It resets counter/cache from every incoming tuple; its resulting top and stack may be arbitrary because the strong block permits them. If $R$ is the successor-based tail encoding of the fixed refutation suffix, the exact first-candidate number is
\[
                         n_y=P(3,P(y,P(1,R)))\ge y.
\]
There is no outer successor on a first candidate. Choose $y\ge\max(1,n)$. Every earlier nonhalting candidate drains, because arithmetic terminates and subsequent positive remainders strictly decrease before tail projection. The scheduler thus reaches a successful candidate unless it halts earlier. No bound on an unknown incoming counter is assumed.
\end{proof}
The verified endpoint has $n=5,d=1$ and otherwise zero theorem data, so this theorem applies directly. No protected-tail debt theorem is needed for this primary construction.
